% vision_performance.tex — Vision Performance Analysis
% Covers: altitude vs pixel size, motion blur, frames on target,
%         coverage completeness, combined detection envelope.
% Insert via % vision_performance.tex — Vision Performance Analysis
% Covers: altitude vs pixel size, motion blur, frames on target,
%         coverage completeness, combined detection envelope.
% Insert via % vision_performance.tex — Vision Performance Analysis
% Covers: altitude vs pixel size, motion blur, frames on target,
%         coverage completeness, combined detection envelope.
% Insert via % vision_performance.tex — Vision Performance Analysis
% Covers: altitude vs pixel size, motion blur, frames on target,
%         coverage completeness, combined detection envelope.
% Insert via \input{sections/vision_performance} in a body or appendix slot.

\section{Vision Performance Analysis}\label{sec:vision-perf}

This section characterises the detection system's operational envelope by
deriving, from first principles, the relationships between altitude, speed,
target pixel size, motion blur, and frame coverage.  All parameters are taken
from the calibrated hardware (Table~\ref{tab:vision-params}); the five
accompanying plots were generated by \texttt{analysis/vision\_performance\_plots.py}
using these values directly.

\begin{table}[ht]
  \centering
  \caption{Camera and inference parameters used for the performance analysis.}\label{tab:vision-params}
  \begin{tabular}{@{}l c l@{}}
    \toprule
    \textbf{Parameter} & \textbf{Value} & \textbf{Source} \\
    \midrule
    Sensor width            & \SI{5.02}{\milli\metre} & IMX296 datasheet \\
    Focal length            & \SI{5.46}{\milli\metre} & Calibrated at \SI{1}{\metre} (2026-03-11) \\
    Sensor resolution       & $1456 \times 1088$\,px  & IMX296 native \\
    Model input             & $640 \times 640$\,px    & YOLOv8n TFLite \\
    Inference time          & \SI{206}{\milli\second}  & Pi~5, XNNPACK, 4 threads \\
    Effective FPS           & 4.8                      & $1000 / 206$ \\
    Shutter                 & 1/1000\,s               & Global shutter (IMX296) \\
    Confidence threshold    & 0.2                      & \texttt{config.py} \\
    Target (dummy)          & \SI{1.8}{\metre} tall    & Mannequin lying on ground \\
    \bottomrule
  \end{tabular}
\end{table}

% ─────────────────────────────────────────────────────────────────────
\subsection{Altitude vs.\ Target Pixel Size}\label{sec:vp:altitude}

The ground footprint width at altitude~$h$ follows from the thin-lens model:
\begin{equation}\label{eq:footprint}
  W_{\mathrm{gnd}} = \frac{w_s \cdot h}{f}\,,
  \qquad
  H_{\mathrm{gnd}} = W_{\mathrm{gnd}} \cdot \frac{N_v}{N_h}
\end{equation}
where $w_s = \SI{5.02}{\milli\metre}$ is the sensor width,
$f = \SI{5.46}{\milli\metre}$ is the focal length, and $N_h \times N_v = 1456
\times 1088$ is the sensor resolution.  The ground sample distance is
$\mathrm{GSD} = W_{\mathrm{gnd}} / N_h$.

The dummy height in the \emph{model input} (640$\times$640) determines
detectability:
\begin{equation}\label{eq:dummypx}
  p_{\mathrm{model}} = \frac{d}{H_{\mathrm{gnd}}} \cdot N_v
    \cdot \frac{640}{\max(N_h, N_v)}
\end{equation}
where $d = \SI{1.8}{\metre}$ is the target height.  Table~\ref{tab:px-by-alt}
lists the result at key altitudes.

\begin{table}[ht]
  \centering
  \caption{Dummy size in the model input at operating altitudes.  The detection
  threshold is approximately \SI{20}{px} in the shortest dimension.}\label{tab:px-by-alt}
  \begin{tabular}{@{}c c c c c@{}}
    \toprule
    \textbf{Altitude (m)} & \textbf{GSD (mm/px)} & \textbf{Full frame (px)} &
      \textbf{Model input (px)} & \textbf{Detectable?} \\
    \midrule
    20 & 12.6 & 143 & 84 & Comfortable \\
    30 & 18.9 &  95 & 56 & Solid \\
    35 & 22.1 &  81 & 48 & Adequate \\
    50 & 31.6 &  57 & 34 & Marginal width \\
    \bottomrule
  \end{tabular}
\end{table}

Figure~\ref{fig:alt-vs-px} shows the continuous relationship.  The dummy
height drops below the \SI{20}{px} YOLOv8 detection threshold at approximately
\SI{85}{\metre}; within the operating range of 20--\SI{50}{\metre} the target
remains at 34--84\,px, comfortably above this floor.  The critical dimension
is the dummy's width (0.5\,m), which shrinks to 7\,px in the model input at
\SI{50}{\metre}---the primary reason for the multi-pass descent strategy
described in Section~\ref{sec:multipass}.

\begin{figure}[ht]
  \centering
  \includegraphics[width=0.85\linewidth]{real_altitude_vs_px}
  \caption{Target height in the 640$\times$640 model input as a function of
  altitude.  The red dashed line marks the 20\,px detection threshold; the
  shaded region indicates unreliable detection.}\label{fig:alt-vs-px}
\end{figure}

% ─────────────────────────────────────────────────────────────────────
\subsection{Speed vs.\ Motion Blur}\label{sec:vp:blur}

The IMX296 is a \emph{global shutter} sensor, eliminating rolling-shutter
artefacts (skew, wobble) entirely.  Translational motion blur remains a
function of exposure time:
\begin{equation}\label{eq:blur}
  b_{\mathrm{px}} = \frac{v \cdot t_{\mathrm{exp}}}{\mathrm{GSD}}
\end{equation}
where $v$ is drone speed and $t_{\mathrm{exp}} = 1/1000$\,s.
Table~\ref{tab:blur} evaluates this across the speed--altitude envelope.

\begin{table}[ht]
  \centering
  \caption{Motion blur in sensor pixels at $1/1000$\,s exposure.  All values
  are sub-pixel at operating speeds.}\label{tab:blur}
  \begin{tabular}{@{}c c c c c@{}}
    \toprule
    \textbf{Speed (m/s)} & \textbf{Exposure} &
      \textbf{20\,m (px)} & \textbf{35\,m (px)} & \textbf{50\,m (px)} \\
    \midrule
     6 & 1/1000\,s & 0.5 & 0.3 & 0.2 \\
    10 & 1/1000\,s & 0.8 & 0.5 & 0.3 \\
    15 & 1/1000\,s & 1.2 & 0.7 & 0.5 \\
    \bottomrule
  \end{tabular}
\end{table}

At the operational speed profile (6--\SI{10}{\metre\per\second}), blur is
0.2--0.8\,px with a $1/1000$\,s shutter---entirely negligible.  Even at
\SI{15}{\metre\per\second} and the lowest altitude of \SI{20}{\metre}, blur
reaches only 1.2\,px, still below the perceptual threshold.
Figure~\ref{fig:speed-blur} plots the continuous relationship at three
altitudes, confirming that the 1\,px visible-blur boundary is only crossed
above \SI{15}{\metre\per\second} at low altitudes.  The global shutter is the
decisive hardware choice: a rolling-shutter sensor would introduce
altitude-dependent skew proportional to the row readout time, adding up to
5--10\,px of geometric distortion at typical multirotor vibration frequencies.

\begin{figure}[ht]
  \centering
  \includegraphics[width=0.85\linewidth]{real_speed_vs_blur}
  \caption{Motion blur vs.\ drone speed at three altitudes (1/1000\,s
  exposure, global shutter).  The grey dashed line marks the 1\,px
  visible-blur threshold.}\label{fig:speed-blur}
\end{figure}

% ─────────────────────────────────────────────────────────────────────
\subsection{Speed vs.\ Frames on Target}\label{sec:vp:frames}

The number of inference frames during which the target is visible determines
detection reliability:
\begin{equation}\label{eq:frames}
  n_{\mathrm{frames}} = \frac{W_{\mathrm{gnd}}}{v} \cdot f_{\mathrm{fps}}
\end{equation}
where $W_{\mathrm{gnd}}$ is the along-track footprint width and
$f_{\mathrm{fps}} = 4.8$.  An adaptive speed profile---linear from
\SI{6}{\metre\per\second} at \SI{20}{\metre} to \SI{10}{\metre\per\second} at
\SI{50}{\metre}---is used during search.

At every altitude in the operating range, the target appears in 11--17 frames
(Table~\ref{tab:frames}), providing substantial margin: a \emph{single}
detection above the 0.2 confidence threshold suffices to trigger
investigation.  The critical speed at which the target would appear in only
one frame is:
\begin{equation}\label{eq:critspeed}
  v_{\mathrm{crit}} = W_{\mathrm{gnd}} \cdot f_{\mathrm{fps}}
\end{equation}
which evaluates to \SI{66}{\metre\per\second} at \SI{20}{\metre} and
\SI{165}{\metre\per\second} at \SI{50}{\metre}---well beyond any multirotor
capability.  Frame coverage is therefore never the limiting factor.

\begin{table}[ht]
  \centering
  \caption{Frames on target at the adaptive speed profile.  Even at maximum
  speed the target is captured in $\geq$11 frames.}\label{tab:frames}
  \begin{tabular}{@{}c c c c c@{}}
    \toprule
    \textbf{Alt.\ (m)} & \textbf{Speed (m/s)} & \textbf{Frame gap (m)} &
      \textbf{Time in view (s)} & \textbf{Frames} \\
    \midrule
    20 &  6.0 & 1.25 & 2.28 & 11.0 \\
    30 &  7.3 & 1.53 & 2.82 & 13.5 \\
    35 &  8.0 & 1.67 & 3.00 & 14.4 \\
    50 & 10.0 & 2.08 & 3.43 & 16.5 \\
    \bottomrule
  \end{tabular}
\end{table}

\begin{figure}[ht]
  \centering
  \includegraphics[width=0.85\linewidth]{real_speed_vs_frames}
  \caption{Frames the target remains visible vs.\ drone speed at three
  altitudes (4.8\,FPS).  The green and red dashed lines mark the 3-frame
  smart-detect threshold and 1-frame minimum, respectively.}\label{fig:speed-frames}
\end{figure}

% ─────────────────────────────────────────────────────────────────────
\subsection{Coverage Efficiency}\label{sec:vp:coverage}

Search efficiency depends on both swath width and speed.  The coverage rate
(area surveyed per unit time) is:
\begin{equation}\label{eq:coverage}
  \dot{A} = H_{\mathrm{gnd}}(h) \times v(h)
\end{equation}
where $H_{\mathrm{gnd}}$ is the cross-track footprint and $v$ is the
altitude-adaptive speed.  At \SI{50}{\metre} the coverage rate is
approximately $343\;\text{m}^2/\text{s}$---roughly 1.7$\times$ that at
\SI{20}{\metre} ($82\;\text{m}^2/\text{s}$).  Starting the search at
\SI{50}{\metre} instead of \SI{35}{\metre} requires approximately 33\% fewer
scan lines and permits 25\% faster flight speed; the combined effect is that
the initial pass completes in roughly 60\% of the time.

The along-track overlap between consecutive frames at altitude~$h$ is:
\begin{equation}\label{eq:overlap}
  \eta = 1 - \frac{v(h)}{W_{\mathrm{gnd}}(h) \cdot f_{\mathrm{fps}}}
\end{equation}
At \SI{35}{\metre} and \SI{8}{\metre\per\second}, $\eta \approx 0.94$ (94\%
frame overlap).  Even at the maximum operating condition (\SI{50}{\metre},
\SI{10}{\metre\per\second}), overlap remains above 93\%.  No part of the
ground passes through the FOV without being imaged in multiple frames.

\begin{figure}[ht]
  \centering
  \includegraphics[width=0.85\linewidth]{real_coverage}
  \caption{Search efficiency vs.\ altitude.  The solid blue line shows
  coverage rate (m$^2$/s); the dashed pink line shows swath width.  Higher
  altitude improves efficiency at the cost of reduced target pixel
  size.}\label{fig:coverage}
\end{figure}

% ─────────────────────────────────────────────────────────────────────
\subsection{Combined Detection Envelope}\label{sec:vp:envelope}

The preceding analyses---pixel size, blur, and frame count---are
combined into a single scalar detection score to identify the operational
boundary:
\begin{equation}\label{eq:score}
  S(v, h) = S_{\mathrm{size}}(h) \times S_{\mathrm{frames}}(v, h)
    \times S_{\mathrm{blur}}(v, h)
\end{equation}
where each component is normalised to $[0, 1]$:
\begin{align}
  S_{\mathrm{size}}   &= \min\!\bigl(p_{\mathrm{model}}(h)\,/\,60,\;1\bigr)
    \label{eq:ssize} \\
  S_{\mathrm{frames}} &= \min\!\bigl(n_{\mathrm{frames}}(v,h)\,/\,5,\;1\bigr)
    \label{eq:sframes} \\
  S_{\mathrm{blur}}   &= 1 - \min\!\bigl(b_{\mathrm{px}}(v,h)\,/\,5,\;1\bigr)
    \label{eq:sblur}
\end{align}
A score of 60\,px in the model input and 5 frames on target represent ``fully
good'' conditions; 5\,px of blur represents total degradation.

Figure~\ref{fig:envelope} renders this score as a speed--altitude heatmap.
The green region ($S > 0.7$) covers the entire operating envelope
(6--\SI{10}{\metre\per\second}, 20--\SI{50}{\metre}), confirming that the
system operates well within its detection limits.  The $S = 0.5$ contour
(marginal detection) is only reached above \SI{55}{\metre} or at speeds
exceeding \SI{12}{\metre\per\second}---neither of which occurs during nominal
operation.

\begin{figure}[ht]
  \centering
  \includegraphics[width=0.88\linewidth]{real_detection_envelope}
  \caption{Combined detection score across the speed--altitude envelope.
  Green indicates high-confidence detection; the white contours mark
  $S = 0.3$, $0.5$, and $0.7$.  The operational regime (6--10\,m/s,
  20--50\,m) lies entirely within the $S > 0.7$ zone.}\label{fig:envelope}
\end{figure}

% ─────────────────────────────────────────────────────────────────────
\subsection{Summary of Limiting Factors}\label{sec:vp:summary}

Table~\ref{tab:limits} summarises the constraint analysis.  The dominant
operational limit is \emph{altitude} (pixel size), not speed or blur.  The
multi-pass descent strategy (50\,m $\to$ 40\,m $\to$ 32\,m) provides a
robust safety net: at 32\,m the dummy occupies 52\,px in the model
input---well above the detection floor.

\begin{table}[ht]
  \centering
  \caption{Summary of vision-system constraints within the operating
  envelope.}\label{tab:limits}
  \begin{tabular}{@{}l c p{6.5cm}@{}}
    \toprule
    \textbf{Factor} & \textbf{Limiting?} & \textbf{Reason} \\
    \midrule
    Motion blur     & No  & Global shutter + 1/1000\,s $\Rightarrow$ sub-pixel blur at all speeds \\
    Pixel size (${<}40$\,m)  & No  & Dummy height $>42$\,px in model input \\
    Pixel size (50\,m)       & Marginal & Dummy height 34\,px, width 7\,px \\
    Inference FPS   & No  & 4.8\,FPS gives 11--17 frames per flyover \\
    Frame gap       & No  & ${\leq}2.1$\,m gap at 10\,m/s; 93\%+ overlap \\
    Speed           & No  & Max 10\,m/s is far below 66\,m/s single-frame limit \\
    \bottomrule
  \end{tabular}
\end{table}

The speed profile could be increased to 15+\,m/s without impacting detection
performance; however, the resulting \SI{1.5}{\metre} GPS timing lag at
\SI{10}{\metre\per\second} (Section~\ref{sec:gps-accuracy}) makes speed
reduction a more effective strategy for improving localisation accuracy than
for preserving detection capability.
 in a body or appendix slot.

\section{Vision Performance Analysis}\label{sec:vision-perf}

This section characterises the detection system's operational envelope by
deriving, from first principles, the relationships between altitude, speed,
target pixel size, motion blur, and frame coverage.  All parameters are taken
from the calibrated hardware (Table~\ref{tab:vision-params}); the five
accompanying plots were generated by \texttt{analysis/vision\_performance\_plots.py}
using these values directly.

\begin{table}[ht]
  \centering
  \caption{Camera and inference parameters used for the performance analysis.}\label{tab:vision-params}
  \begin{tabular}{@{}l c l@{}}
    \toprule
    \textbf{Parameter} & \textbf{Value} & \textbf{Source} \\
    \midrule
    Sensor width            & \SI{5.02}{\milli\metre} & IMX296 datasheet \\
    Focal length            & \SI{5.46}{\milli\metre} & Calibrated at \SI{1}{\metre} (2026-03-11) \\
    Sensor resolution       & $1456 \times 1088$\,px  & IMX296 native \\
    Model input             & $640 \times 640$\,px    & YOLOv8n TFLite \\
    Inference time          & \SI{206}{\milli\second}  & Pi~5, XNNPACK, 4 threads \\
    Effective FPS           & 4.8                      & $1000 / 206$ \\
    Shutter                 & 1/1000\,s               & Global shutter (IMX296) \\
    Confidence threshold    & 0.2                      & \texttt{config.py} \\
    Target (dummy)          & \SI{1.8}{\metre} tall    & Mannequin lying on ground \\
    \bottomrule
  \end{tabular}
\end{table}

% ─────────────────────────────────────────────────────────────────────
\subsection{Altitude vs.\ Target Pixel Size}\label{sec:vp:altitude}

The ground footprint width at altitude~$h$ follows from the thin-lens model:
\begin{equation}\label{eq:footprint}
  W_{\mathrm{gnd}} = \frac{w_s \cdot h}{f}\,,
  \qquad
  H_{\mathrm{gnd}} = W_{\mathrm{gnd}} \cdot \frac{N_v}{N_h}
\end{equation}
where $w_s = \SI{5.02}{\milli\metre}$ is the sensor width,
$f = \SI{5.46}{\milli\metre}$ is the focal length, and $N_h \times N_v = 1456
\times 1088$ is the sensor resolution.  The ground sample distance is
$\mathrm{GSD} = W_{\mathrm{gnd}} / N_h$.

The dummy height in the \emph{model input} (640$\times$640) determines
detectability:
\begin{equation}\label{eq:dummypx}
  p_{\mathrm{model}} = \frac{d}{H_{\mathrm{gnd}}} \cdot N_v
    \cdot \frac{640}{\max(N_h, N_v)}
\end{equation}
where $d = \SI{1.8}{\metre}$ is the target height.  Table~\ref{tab:px-by-alt}
lists the result at key altitudes.

\begin{table}[ht]
  \centering
  \caption{Dummy size in the model input at operating altitudes.  The detection
  threshold is approximately \SI{20}{px} in the shortest dimension.}\label{tab:px-by-alt}
  \begin{tabular}{@{}c c c c c@{}}
    \toprule
    \textbf{Altitude (m)} & \textbf{GSD (mm/px)} & \textbf{Full frame (px)} &
      \textbf{Model input (px)} & \textbf{Detectable?} \\
    \midrule
    20 & 12.6 & 143 & 84 & Comfortable \\
    30 & 18.9 &  95 & 56 & Solid \\
    35 & 22.1 &  81 & 48 & Adequate \\
    50 & 31.6 &  57 & 34 & Marginal width \\
    \bottomrule
  \end{tabular}
\end{table}

Figure~\ref{fig:alt-vs-px} shows the continuous relationship.  The dummy
height drops below the \SI{20}{px} YOLOv8 detection threshold at approximately
\SI{85}{\metre}; within the operating range of 20--\SI{50}{\metre} the target
remains at 34--84\,px, comfortably above this floor.  The critical dimension
is the dummy's width (0.5\,m), which shrinks to 7\,px in the model input at
\SI{50}{\metre}---the primary reason for the multi-pass descent strategy
described in Section~\ref{sec:multipass}.

\begin{figure}[ht]
  \centering
  \includegraphics[width=0.85\linewidth]{real_altitude_vs_px}
  \caption{Target height in the 640$\times$640 model input as a function of
  altitude.  The red dashed line marks the 20\,px detection threshold; the
  shaded region indicates unreliable detection.}\label{fig:alt-vs-px}
\end{figure}

% ─────────────────────────────────────────────────────────────────────
\subsection{Speed vs.\ Motion Blur}\label{sec:vp:blur}

The IMX296 is a \emph{global shutter} sensor, eliminating rolling-shutter
artefacts (skew, wobble) entirely.  Translational motion blur remains a
function of exposure time:
\begin{equation}\label{eq:blur}
  b_{\mathrm{px}} = \frac{v \cdot t_{\mathrm{exp}}}{\mathrm{GSD}}
\end{equation}
where $v$ is drone speed and $t_{\mathrm{exp}} = 1/1000$\,s.
Table~\ref{tab:blur} evaluates this across the speed--altitude envelope.

\begin{table}[ht]
  \centering
  \caption{Motion blur in sensor pixels at $1/1000$\,s exposure.  All values
  are sub-pixel at operating speeds.}\label{tab:blur}
  \begin{tabular}{@{}c c c c c@{}}
    \toprule
    \textbf{Speed (m/s)} & \textbf{Exposure} &
      \textbf{20\,m (px)} & \textbf{35\,m (px)} & \textbf{50\,m (px)} \\
    \midrule
     6 & 1/1000\,s & 0.5 & 0.3 & 0.2 \\
    10 & 1/1000\,s & 0.8 & 0.5 & 0.3 \\
    15 & 1/1000\,s & 1.2 & 0.7 & 0.5 \\
    \bottomrule
  \end{tabular}
\end{table}

At the operational speed profile (6--\SI{10}{\metre\per\second}), blur is
0.2--0.8\,px with a $1/1000$\,s shutter---entirely negligible.  Even at
\SI{15}{\metre\per\second} and the lowest altitude of \SI{20}{\metre}, blur
reaches only 1.2\,px, still below the perceptual threshold.
Figure~\ref{fig:speed-blur} plots the continuous relationship at three
altitudes, confirming that the 1\,px visible-blur boundary is only crossed
above \SI{15}{\metre\per\second} at low altitudes.  The global shutter is the
decisive hardware choice: a rolling-shutter sensor would introduce
altitude-dependent skew proportional to the row readout time, adding up to
5--10\,px of geometric distortion at typical multirotor vibration frequencies.

\begin{figure}[ht]
  \centering
  \includegraphics[width=0.85\linewidth]{real_speed_vs_blur}
  \caption{Motion blur vs.\ drone speed at three altitudes (1/1000\,s
  exposure, global shutter).  The grey dashed line marks the 1\,px
  visible-blur threshold.}\label{fig:speed-blur}
\end{figure}

% ─────────────────────────────────────────────────────────────────────
\subsection{Speed vs.\ Frames on Target}\label{sec:vp:frames}

The number of inference frames during which the target is visible determines
detection reliability:
\begin{equation}\label{eq:frames}
  n_{\mathrm{frames}} = \frac{W_{\mathrm{gnd}}}{v} \cdot f_{\mathrm{fps}}
\end{equation}
where $W_{\mathrm{gnd}}$ is the along-track footprint width and
$f_{\mathrm{fps}} = 4.8$.  An adaptive speed profile---linear from
\SI{6}{\metre\per\second} at \SI{20}{\metre} to \SI{10}{\metre\per\second} at
\SI{50}{\metre}---is used during search.

At every altitude in the operating range, the target appears in 11--17 frames
(Table~\ref{tab:frames}), providing substantial margin: a \emph{single}
detection above the 0.2 confidence threshold suffices to trigger
investigation.  The critical speed at which the target would appear in only
one frame is:
\begin{equation}\label{eq:critspeed}
  v_{\mathrm{crit}} = W_{\mathrm{gnd}} \cdot f_{\mathrm{fps}}
\end{equation}
which evaluates to \SI{66}{\metre\per\second} at \SI{20}{\metre} and
\SI{165}{\metre\per\second} at \SI{50}{\metre}---well beyond any multirotor
capability.  Frame coverage is therefore never the limiting factor.

\begin{table}[ht]
  \centering
  \caption{Frames on target at the adaptive speed profile.  Even at maximum
  speed the target is captured in $\geq$11 frames.}\label{tab:frames}
  \begin{tabular}{@{}c c c c c@{}}
    \toprule
    \textbf{Alt.\ (m)} & \textbf{Speed (m/s)} & \textbf{Frame gap (m)} &
      \textbf{Time in view (s)} & \textbf{Frames} \\
    \midrule
    20 &  6.0 & 1.25 & 2.28 & 11.0 \\
    30 &  7.3 & 1.53 & 2.82 & 13.5 \\
    35 &  8.0 & 1.67 & 3.00 & 14.4 \\
    50 & 10.0 & 2.08 & 3.43 & 16.5 \\
    \bottomrule
  \end{tabular}
\end{table}

\begin{figure}[ht]
  \centering
  \includegraphics[width=0.85\linewidth]{real_speed_vs_frames}
  \caption{Frames the target remains visible vs.\ drone speed at three
  altitudes (4.8\,FPS).  The green and red dashed lines mark the 3-frame
  smart-detect threshold and 1-frame minimum, respectively.}\label{fig:speed-frames}
\end{figure}

% ─────────────────────────────────────────────────────────────────────
\subsection{Coverage Efficiency}\label{sec:vp:coverage}

Search efficiency depends on both swath width and speed.  The coverage rate
(area surveyed per unit time) is:
\begin{equation}\label{eq:coverage}
  \dot{A} = H_{\mathrm{gnd}}(h) \times v(h)
\end{equation}
where $H_{\mathrm{gnd}}$ is the cross-track footprint and $v$ is the
altitude-adaptive speed.  At \SI{50}{\metre} the coverage rate is
approximately $343\;\text{m}^2/\text{s}$---roughly 1.7$\times$ that at
\SI{20}{\metre} ($82\;\text{m}^2/\text{s}$).  Starting the search at
\SI{50}{\metre} instead of \SI{35}{\metre} requires approximately 33\% fewer
scan lines and permits 25\% faster flight speed; the combined effect is that
the initial pass completes in roughly 60\% of the time.

The along-track overlap between consecutive frames at altitude~$h$ is:
\begin{equation}\label{eq:overlap}
  \eta = 1 - \frac{v(h)}{W_{\mathrm{gnd}}(h) \cdot f_{\mathrm{fps}}}
\end{equation}
At \SI{35}{\metre} and \SI{8}{\metre\per\second}, $\eta \approx 0.94$ (94\%
frame overlap).  Even at the maximum operating condition (\SI{50}{\metre},
\SI{10}{\metre\per\second}), overlap remains above 93\%.  No part of the
ground passes through the FOV without being imaged in multiple frames.

\begin{figure}[ht]
  \centering
  \includegraphics[width=0.85\linewidth]{real_coverage}
  \caption{Search efficiency vs.\ altitude.  The solid blue line shows
  coverage rate (m$^2$/s); the dashed pink line shows swath width.  Higher
  altitude improves efficiency at the cost of reduced target pixel
  size.}\label{fig:coverage}
\end{figure}

% ─────────────────────────────────────────────────────────────────────
\subsection{Combined Detection Envelope}\label{sec:vp:envelope}

The preceding analyses---pixel size, blur, and frame count---are
combined into a single scalar detection score to identify the operational
boundary:
\begin{equation}\label{eq:score}
  S(v, h) = S_{\mathrm{size}}(h) \times S_{\mathrm{frames}}(v, h)
    \times S_{\mathrm{blur}}(v, h)
\end{equation}
where each component is normalised to $[0, 1]$:
\begin{align}
  S_{\mathrm{size}}   &= \min\!\bigl(p_{\mathrm{model}}(h)\,/\,60,\;1\bigr)
    \label{eq:ssize} \\
  S_{\mathrm{frames}} &= \min\!\bigl(n_{\mathrm{frames}}(v,h)\,/\,5,\;1\bigr)
    \label{eq:sframes} \\
  S_{\mathrm{blur}}   &= 1 - \min\!\bigl(b_{\mathrm{px}}(v,h)\,/\,5,\;1\bigr)
    \label{eq:sblur}
\end{align}
A score of 60\,px in the model input and 5 frames on target represent ``fully
good'' conditions; 5\,px of blur represents total degradation.

Figure~\ref{fig:envelope} renders this score as a speed--altitude heatmap.
The green region ($S > 0.7$) covers the entire operating envelope
(6--\SI{10}{\metre\per\second}, 20--\SI{50}{\metre}), confirming that the
system operates well within its detection limits.  The $S = 0.5$ contour
(marginal detection) is only reached above \SI{55}{\metre} or at speeds
exceeding \SI{12}{\metre\per\second}---neither of which occurs during nominal
operation.

\begin{figure}[ht]
  \centering
  \includegraphics[width=0.88\linewidth]{real_detection_envelope}
  \caption{Combined detection score across the speed--altitude envelope.
  Green indicates high-confidence detection; the white contours mark
  $S = 0.3$, $0.5$, and $0.7$.  The operational regime (6--10\,m/s,
  20--50\,m) lies entirely within the $S > 0.7$ zone.}\label{fig:envelope}
\end{figure}

% ─────────────────────────────────────────────────────────────────────
\subsection{Summary of Limiting Factors}\label{sec:vp:summary}

Table~\ref{tab:limits} summarises the constraint analysis.  The dominant
operational limit is \emph{altitude} (pixel size), not speed or blur.  The
multi-pass descent strategy (50\,m $\to$ 40\,m $\to$ 32\,m) provides a
robust safety net: at 32\,m the dummy occupies 52\,px in the model
input---well above the detection floor.

\begin{table}[ht]
  \centering
  \caption{Summary of vision-system constraints within the operating
  envelope.}\label{tab:limits}
  \begin{tabular}{@{}l c p{6.5cm}@{}}
    \toprule
    \textbf{Factor} & \textbf{Limiting?} & \textbf{Reason} \\
    \midrule
    Motion blur     & No  & Global shutter + 1/1000\,s $\Rightarrow$ sub-pixel blur at all speeds \\
    Pixel size (${<}40$\,m)  & No  & Dummy height $>42$\,px in model input \\
    Pixel size (50\,m)       & Marginal & Dummy height 34\,px, width 7\,px \\
    Inference FPS   & No  & 4.8\,FPS gives 11--17 frames per flyover \\
    Frame gap       & No  & ${\leq}2.1$\,m gap at 10\,m/s; 93\%+ overlap \\
    Speed           & No  & Max 10\,m/s is far below 66\,m/s single-frame limit \\
    \bottomrule
  \end{tabular}
\end{table}

The speed profile could be increased to 15+\,m/s without impacting detection
performance; however, the resulting \SI{1.5}{\metre} GPS timing lag at
\SI{10}{\metre\per\second} (Section~\ref{sec:gps-accuracy}) makes speed
reduction a more effective strategy for improving localisation accuracy than
for preserving detection capability.
 in a body or appendix slot.

\section{Vision Performance Analysis}\label{sec:vision-perf}

This section characterises the detection system's operational envelope by
deriving, from first principles, the relationships between altitude, speed,
target pixel size, motion blur, and frame coverage.  All parameters are taken
from the calibrated hardware (Table~\ref{tab:vision-params}); the five
accompanying plots were generated by \texttt{analysis/vision\_performance\_plots.py}
using these values directly.

\begin{table}[ht]
  \centering
  \caption{Camera and inference parameters used for the performance analysis.}\label{tab:vision-params}
  \begin{tabular}{@{}l c l@{}}
    \toprule
    \textbf{Parameter} & \textbf{Value} & \textbf{Source} \\
    \midrule
    Sensor width            & \SI{5.02}{\milli\metre} & IMX296 datasheet \\
    Focal length            & \SI{5.46}{\milli\metre} & Calibrated at \SI{1}{\metre} (2026-03-11) \\
    Sensor resolution       & $1456 \times 1088$\,px  & IMX296 native \\
    Model input             & $640 \times 640$\,px    & YOLOv8n TFLite \\
    Inference time          & \SI{206}{\milli\second}  & Pi~5, XNNPACK, 4 threads \\
    Effective FPS           & 4.8                      & $1000 / 206$ \\
    Shutter                 & 1/1000\,s               & Global shutter (IMX296) \\
    Confidence threshold    & 0.2                      & \texttt{config.py} \\
    Target (dummy)          & \SI{1.8}{\metre} tall    & Mannequin lying on ground \\
    \bottomrule
  \end{tabular}
\end{table}

% ─────────────────────────────────────────────────────────────────────
\subsection{Altitude vs.\ Target Pixel Size}\label{sec:vp:altitude}

The ground footprint width at altitude~$h$ follows from the thin-lens model:
\begin{equation}\label{eq:footprint}
  W_{\mathrm{gnd}} = \frac{w_s \cdot h}{f}\,,
  \qquad
  H_{\mathrm{gnd}} = W_{\mathrm{gnd}} \cdot \frac{N_v}{N_h}
\end{equation}
where $w_s = \SI{5.02}{\milli\metre}$ is the sensor width,
$f = \SI{5.46}{\milli\metre}$ is the focal length, and $N_h \times N_v = 1456
\times 1088$ is the sensor resolution.  The ground sample distance is
$\mathrm{GSD} = W_{\mathrm{gnd}} / N_h$.

The dummy height in the \emph{model input} (640$\times$640) determines
detectability:
\begin{equation}\label{eq:dummypx}
  p_{\mathrm{model}} = \frac{d}{H_{\mathrm{gnd}}} \cdot N_v
    \cdot \frac{640}{\max(N_h, N_v)}
\end{equation}
where $d = \SI{1.8}{\metre}$ is the target height.  Table~\ref{tab:px-by-alt}
lists the result at key altitudes.

\begin{table}[ht]
  \centering
  \caption{Dummy size in the model input at operating altitudes.  The detection
  threshold is approximately \SI{20}{px} in the shortest dimension.}\label{tab:px-by-alt}
  \begin{tabular}{@{}c c c c c@{}}
    \toprule
    \textbf{Altitude (m)} & \textbf{GSD (mm/px)} & \textbf{Full frame (px)} &
      \textbf{Model input (px)} & \textbf{Detectable?} \\
    \midrule
    20 & 12.6 & 143 & 84 & Comfortable \\
    30 & 18.9 &  95 & 56 & Solid \\
    35 & 22.1 &  81 & 48 & Adequate \\
    50 & 31.6 &  57 & 34 & Marginal width \\
    \bottomrule
  \end{tabular}
\end{table}

Figure~\ref{fig:alt-vs-px} shows the continuous relationship.  The dummy
height drops below the \SI{20}{px} YOLOv8 detection threshold at approximately
\SI{85}{\metre}; within the operating range of 20--\SI{50}{\metre} the target
remains at 34--84\,px, comfortably above this floor.  The critical dimension
is the dummy's width (0.5\,m), which shrinks to 7\,px in the model input at
\SI{50}{\metre}---the primary reason for the multi-pass descent strategy
described in Section~\ref{sec:multipass}.

\begin{figure}[ht]
  \centering
  \includegraphics[width=0.85\linewidth]{real_altitude_vs_px}
  \caption{Target height in the 640$\times$640 model input as a function of
  altitude.  The red dashed line marks the 20\,px detection threshold; the
  shaded region indicates unreliable detection.}\label{fig:alt-vs-px}
\end{figure}

% ─────────────────────────────────────────────────────────────────────
\subsection{Speed vs.\ Motion Blur}\label{sec:vp:blur}

The IMX296 is a \emph{global shutter} sensor, eliminating rolling-shutter
artefacts (skew, wobble) entirely.  Translational motion blur remains a
function of exposure time:
\begin{equation}\label{eq:blur}
  b_{\mathrm{px}} = \frac{v \cdot t_{\mathrm{exp}}}{\mathrm{GSD}}
\end{equation}
where $v$ is drone speed and $t_{\mathrm{exp}} = 1/1000$\,s.
Table~\ref{tab:blur} evaluates this across the speed--altitude envelope.

\begin{table}[ht]
  \centering
  \caption{Motion blur in sensor pixels at $1/1000$\,s exposure.  All values
  are sub-pixel at operating speeds.}\label{tab:blur}
  \begin{tabular}{@{}c c c c c@{}}
    \toprule
    \textbf{Speed (m/s)} & \textbf{Exposure} &
      \textbf{20\,m (px)} & \textbf{35\,m (px)} & \textbf{50\,m (px)} \\
    \midrule
     6 & 1/1000\,s & 0.5 & 0.3 & 0.2 \\
    10 & 1/1000\,s & 0.8 & 0.5 & 0.3 \\
    15 & 1/1000\,s & 1.2 & 0.7 & 0.5 \\
    \bottomrule
  \end{tabular}
\end{table}

At the operational speed profile (6--\SI{10}{\metre\per\second}), blur is
0.2--0.8\,px with a $1/1000$\,s shutter---entirely negligible.  Even at
\SI{15}{\metre\per\second} and the lowest altitude of \SI{20}{\metre}, blur
reaches only 1.2\,px, still below the perceptual threshold.
Figure~\ref{fig:speed-blur} plots the continuous relationship at three
altitudes, confirming that the 1\,px visible-blur boundary is only crossed
above \SI{15}{\metre\per\second} at low altitudes.  The global shutter is the
decisive hardware choice: a rolling-shutter sensor would introduce
altitude-dependent skew proportional to the row readout time, adding up to
5--10\,px of geometric distortion at typical multirotor vibration frequencies.

\begin{figure}[ht]
  \centering
  \includegraphics[width=0.85\linewidth]{real_speed_vs_blur}
  \caption{Motion blur vs.\ drone speed at three altitudes (1/1000\,s
  exposure, global shutter).  The grey dashed line marks the 1\,px
  visible-blur threshold.}\label{fig:speed-blur}
\end{figure}

% ─────────────────────────────────────────────────────────────────────
\subsection{Speed vs.\ Frames on Target}\label{sec:vp:frames}

The number of inference frames during which the target is visible determines
detection reliability:
\begin{equation}\label{eq:frames}
  n_{\mathrm{frames}} = \frac{W_{\mathrm{gnd}}}{v} \cdot f_{\mathrm{fps}}
\end{equation}
where $W_{\mathrm{gnd}}$ is the along-track footprint width and
$f_{\mathrm{fps}} = 4.8$.  An adaptive speed profile---linear from
\SI{6}{\metre\per\second} at \SI{20}{\metre} to \SI{10}{\metre\per\second} at
\SI{50}{\metre}---is used during search.

At every altitude in the operating range, the target appears in 11--17 frames
(Table~\ref{tab:frames}), providing substantial margin: a \emph{single}
detection above the 0.2 confidence threshold suffices to trigger
investigation.  The critical speed at which the target would appear in only
one frame is:
\begin{equation}\label{eq:critspeed}
  v_{\mathrm{crit}} = W_{\mathrm{gnd}} \cdot f_{\mathrm{fps}}
\end{equation}
which evaluates to \SI{66}{\metre\per\second} at \SI{20}{\metre} and
\SI{165}{\metre\per\second} at \SI{50}{\metre}---well beyond any multirotor
capability.  Frame coverage is therefore never the limiting factor.

\begin{table}[ht]
  \centering
  \caption{Frames on target at the adaptive speed profile.  Even at maximum
  speed the target is captured in $\geq$11 frames.}\label{tab:frames}
  \begin{tabular}{@{}c c c c c@{}}
    \toprule
    \textbf{Alt.\ (m)} & \textbf{Speed (m/s)} & \textbf{Frame gap (m)} &
      \textbf{Time in view (s)} & \textbf{Frames} \\
    \midrule
    20 &  6.0 & 1.25 & 2.28 & 11.0 \\
    30 &  7.3 & 1.53 & 2.82 & 13.5 \\
    35 &  8.0 & 1.67 & 3.00 & 14.4 \\
    50 & 10.0 & 2.08 & 3.43 & 16.5 \\
    \bottomrule
  \end{tabular}
\end{table}

\begin{figure}[ht]
  \centering
  \includegraphics[width=0.85\linewidth]{real_speed_vs_frames}
  \caption{Frames the target remains visible vs.\ drone speed at three
  altitudes (4.8\,FPS).  The green and red dashed lines mark the 3-frame
  smart-detect threshold and 1-frame minimum, respectively.}\label{fig:speed-frames}
\end{figure}

% ─────────────────────────────────────────────────────────────────────
\subsection{Coverage Efficiency}\label{sec:vp:coverage}

Search efficiency depends on both swath width and speed.  The coverage rate
(area surveyed per unit time) is:
\begin{equation}\label{eq:coverage}
  \dot{A} = H_{\mathrm{gnd}}(h) \times v(h)
\end{equation}
where $H_{\mathrm{gnd}}$ is the cross-track footprint and $v$ is the
altitude-adaptive speed.  At \SI{50}{\metre} the coverage rate is
approximately $343\;\text{m}^2/\text{s}$---roughly 1.7$\times$ that at
\SI{20}{\metre} ($82\;\text{m}^2/\text{s}$).  Starting the search at
\SI{50}{\metre} instead of \SI{35}{\metre} requires approximately 33\% fewer
scan lines and permits 25\% faster flight speed; the combined effect is that
the initial pass completes in roughly 60\% of the time.

The along-track overlap between consecutive frames at altitude~$h$ is:
\begin{equation}\label{eq:overlap}
  \eta = 1 - \frac{v(h)}{W_{\mathrm{gnd}}(h) \cdot f_{\mathrm{fps}}}
\end{equation}
At \SI{35}{\metre} and \SI{8}{\metre\per\second}, $\eta \approx 0.94$ (94\%
frame overlap).  Even at the maximum operating condition (\SI{50}{\metre},
\SI{10}{\metre\per\second}), overlap remains above 93\%.  No part of the
ground passes through the FOV without being imaged in multiple frames.

\begin{figure}[ht]
  \centering
  \includegraphics[width=0.85\linewidth]{real_coverage}
  \caption{Search efficiency vs.\ altitude.  The solid blue line shows
  coverage rate (m$^2$/s); the dashed pink line shows swath width.  Higher
  altitude improves efficiency at the cost of reduced target pixel
  size.}\label{fig:coverage}
\end{figure}

% ─────────────────────────────────────────────────────────────────────
\subsection{Combined Detection Envelope}\label{sec:vp:envelope}

The preceding analyses---pixel size, blur, and frame count---are
combined into a single scalar detection score to identify the operational
boundary:
\begin{equation}\label{eq:score}
  S(v, h) = S_{\mathrm{size}}(h) \times S_{\mathrm{frames}}(v, h)
    \times S_{\mathrm{blur}}(v, h)
\end{equation}
where each component is normalised to $[0, 1]$:
\begin{align}
  S_{\mathrm{size}}   &= \min\!\bigl(p_{\mathrm{model}}(h)\,/\,60,\;1\bigr)
    \label{eq:ssize} \\
  S_{\mathrm{frames}} &= \min\!\bigl(n_{\mathrm{frames}}(v,h)\,/\,5,\;1\bigr)
    \label{eq:sframes} \\
  S_{\mathrm{blur}}   &= 1 - \min\!\bigl(b_{\mathrm{px}}(v,h)\,/\,5,\;1\bigr)
    \label{eq:sblur}
\end{align}
A score of 60\,px in the model input and 5 frames on target represent ``fully
good'' conditions; 5\,px of blur represents total degradation.

Figure~\ref{fig:envelope} renders this score as a speed--altitude heatmap.
The green region ($S > 0.7$) covers the entire operating envelope
(6--\SI{10}{\metre\per\second}, 20--\SI{50}{\metre}), confirming that the
system operates well within its detection limits.  The $S = 0.5$ contour
(marginal detection) is only reached above \SI{55}{\metre} or at speeds
exceeding \SI{12}{\metre\per\second}---neither of which occurs during nominal
operation.

\begin{figure}[ht]
  \centering
  \includegraphics[width=0.88\linewidth]{real_detection_envelope}
  \caption{Combined detection score across the speed--altitude envelope.
  Green indicates high-confidence detection; the white contours mark
  $S = 0.3$, $0.5$, and $0.7$.  The operational regime (6--10\,m/s,
  20--50\,m) lies entirely within the $S > 0.7$ zone.}\label{fig:envelope}
\end{figure}

% ─────────────────────────────────────────────────────────────────────
\subsection{Summary of Limiting Factors}\label{sec:vp:summary}

Table~\ref{tab:limits} summarises the constraint analysis.  The dominant
operational limit is \emph{altitude} (pixel size), not speed or blur.  The
multi-pass descent strategy (50\,m $\to$ 40\,m $\to$ 32\,m) provides a
robust safety net: at 32\,m the dummy occupies 52\,px in the model
input---well above the detection floor.

\begin{table}[ht]
  \centering
  \caption{Summary of vision-system constraints within the operating
  envelope.}\label{tab:limits}
  \begin{tabular}{@{}l c p{6.5cm}@{}}
    \toprule
    \textbf{Factor} & \textbf{Limiting?} & \textbf{Reason} \\
    \midrule
    Motion blur     & No  & Global shutter + 1/1000\,s $\Rightarrow$ sub-pixel blur at all speeds \\
    Pixel size (${<}40$\,m)  & No  & Dummy height $>42$\,px in model input \\
    Pixel size (50\,m)       & Marginal & Dummy height 34\,px, width 7\,px \\
    Inference FPS   & No  & 4.8\,FPS gives 11--17 frames per flyover \\
    Frame gap       & No  & ${\leq}2.1$\,m gap at 10\,m/s; 93\%+ overlap \\
    Speed           & No  & Max 10\,m/s is far below 66\,m/s single-frame limit \\
    \bottomrule
  \end{tabular}
\end{table}

The speed profile could be increased to 15+\,m/s without impacting detection
performance; however, the resulting \SI{1.5}{\metre} GPS timing lag at
\SI{10}{\metre\per\second} (Section~\ref{sec:gps-accuracy}) makes speed
reduction a more effective strategy for improving localisation accuracy than
for preserving detection capability.
 in a body or appendix slot.

\section{Vision Performance Analysis}\label{sec:vision-perf}

This section characterises the detection system's operational envelope by
deriving, from first principles, the relationships between altitude, speed,
target pixel size, motion blur, and frame coverage.  All parameters are taken
from the calibrated hardware (Table~\ref{tab:vision-params}); the five
accompanying plots were generated by \texttt{analysis/vision\_performance\_plots.py}
using these values directly.

\begin{table}[ht]
  \centering
  \caption{Camera and inference parameters used for the performance analysis.}\label{tab:vision-params}
  \begin{tabular}{@{}l c l@{}}
    \toprule
    \textbf{Parameter} & \textbf{Value} & \textbf{Source} \\
    \midrule
    Sensor width            & \SI{5.02}{\milli\metre} & IMX296 datasheet \\
    Focal length            & \SI{5.46}{\milli\metre} & Calibrated at \SI{1}{\metre} (2026-03-11) \\
    Sensor resolution       & $1456 \times 1088$\,px  & IMX296 native \\
    Model input             & $640 \times 640$\,px    & YOLOv8n TFLite \\
    Inference time          & \SI{206}{\milli\second}  & Pi~5, XNNPACK, 4 threads \\
    Effective FPS           & 4.8                      & $1000 / 206$ \\
    Shutter                 & 1/1000\,s               & Global shutter (IMX296) \\
    Confidence threshold    & 0.2                      & \texttt{config.py} \\
    Target (dummy)          & \SI{1.8}{\metre} tall    & Mannequin lying on ground \\
    \bottomrule
  \end{tabular}
\end{table}

% ─────────────────────────────────────────────────────────────────────
\subsection{Altitude vs.\ Target Pixel Size}\label{sec:vp:altitude}

The ground footprint width at altitude~$h$ follows from the thin-lens model:
\begin{equation}\label{eq:footprint}
  W_{\mathrm{gnd}} = \frac{w_s \cdot h}{f}\,,
  \qquad
  H_{\mathrm{gnd}} = W_{\mathrm{gnd}} \cdot \frac{N_v}{N_h}
\end{equation}
where $w_s = \SI{5.02}{\milli\metre}$ is the sensor width,
$f = \SI{5.46}{\milli\metre}$ is the focal length, and $N_h \times N_v = 1456
\times 1088$ is the sensor resolution.  The ground sample distance is
$\mathrm{GSD} = W_{\mathrm{gnd}} / N_h$.

The dummy height in the \emph{model input} (640$\times$640) determines
detectability:
\begin{equation}\label{eq:dummypx}
  p_{\mathrm{model}} = \frac{d}{H_{\mathrm{gnd}}} \cdot N_v
    \cdot \frac{640}{\max(N_h, N_v)}
\end{equation}
where $d = \SI{1.8}{\metre}$ is the target height.  Table~\ref{tab:px-by-alt}
lists the result at key altitudes.

\begin{table}[ht]
  \centering
  \caption{Dummy size in the model input at operating altitudes.  The detection
  threshold is approximately \SI{20}{px} in the shortest dimension.}\label{tab:px-by-alt}
  \begin{tabular}{@{}c c c c c@{}}
    \toprule
    \textbf{Altitude (m)} & \textbf{GSD (mm/px)} & \textbf{Full frame (px)} &
      \textbf{Model input (px)} & \textbf{Detectable?} \\
    \midrule
    20 & 12.6 & 143 & 84 & Comfortable \\
    30 & 18.9 &  95 & 56 & Solid \\
    35 & 22.1 &  81 & 48 & Adequate \\
    50 & 31.6 &  57 & 34 & Marginal width \\
    \bottomrule
  \end{tabular}
\end{table}

Figure~\ref{fig:alt-vs-px} shows the continuous relationship.  The dummy
height drops below the \SI{20}{px} YOLOv8 detection threshold at approximately
\SI{85}{\metre}; within the operating range of 20--\SI{50}{\metre} the target
remains at 34--84\,px, comfortably above this floor.  The critical dimension
is the dummy's width (0.5\,m), which shrinks to 7\,px in the model input at
\SI{50}{\metre}---the primary reason for the multi-pass descent strategy
described in Section~\ref{sec:multipass}.

\begin{figure}[ht]
  \centering
  \includegraphics[width=0.85\linewidth]{real_altitude_vs_px}
  \caption{Target height in the 640$\times$640 model input as a function of
  altitude.  The red dashed line marks the 20\,px detection threshold; the
  shaded region indicates unreliable detection.}\label{fig:alt-vs-px}
\end{figure}

% ─────────────────────────────────────────────────────────────────────
\subsection{Speed vs.\ Motion Blur}\label{sec:vp:blur}

The IMX296 is a \emph{global shutter} sensor, eliminating rolling-shutter
artefacts (skew, wobble) entirely.  Translational motion blur remains a
function of exposure time:
\begin{equation}\label{eq:blur}
  b_{\mathrm{px}} = \frac{v \cdot t_{\mathrm{exp}}}{\mathrm{GSD}}
\end{equation}
where $v$ is drone speed and $t_{\mathrm{exp}} = 1/1000$\,s.
Table~\ref{tab:blur} evaluates this across the speed--altitude envelope.

\begin{table}[ht]
  \centering
  \caption{Motion blur in sensor pixels at $1/1000$\,s exposure.  All values
  are sub-pixel at operating speeds.}\label{tab:blur}
  \begin{tabular}{@{}c c c c c@{}}
    \toprule
    \textbf{Speed (m/s)} & \textbf{Exposure} &
      \textbf{20\,m (px)} & \textbf{35\,m (px)} & \textbf{50\,m (px)} \\
    \midrule
     6 & 1/1000\,s & 0.5 & 0.3 & 0.2 \\
    10 & 1/1000\,s & 0.8 & 0.5 & 0.3 \\
    15 & 1/1000\,s & 1.2 & 0.7 & 0.5 \\
    \bottomrule
  \end{tabular}
\end{table}

At the operational speed profile (6--\SI{10}{\metre\per\second}), blur is
0.2--0.8\,px with a $1/1000$\,s shutter---entirely negligible.  Even at
\SI{15}{\metre\per\second} and the lowest altitude of \SI{20}{\metre}, blur
reaches only 1.2\,px, still below the perceptual threshold.
Figure~\ref{fig:speed-blur} plots the continuous relationship at three
altitudes, confirming that the 1\,px visible-blur boundary is only crossed
above \SI{15}{\metre\per\second} at low altitudes.  The global shutter is the
decisive hardware choice: a rolling-shutter sensor would introduce
altitude-dependent skew proportional to the row readout time, adding up to
5--10\,px of geometric distortion at typical multirotor vibration frequencies.

\begin{figure}[ht]
  \centering
  \includegraphics[width=0.85\linewidth]{real_speed_vs_blur}
  \caption{Motion blur vs.\ drone speed at three altitudes (1/1000\,s
  exposure, global shutter).  The grey dashed line marks the 1\,px
  visible-blur threshold.}\label{fig:speed-blur}
\end{figure}

% ─────────────────────────────────────────────────────────────────────
\subsection{Speed vs.\ Frames on Target}\label{sec:vp:frames}

The number of inference frames during which the target is visible determines
detection reliability:
\begin{equation}\label{eq:frames}
  n_{\mathrm{frames}} = \frac{W_{\mathrm{gnd}}}{v} \cdot f_{\mathrm{fps}}
\end{equation}
where $W_{\mathrm{gnd}}$ is the along-track footprint width and
$f_{\mathrm{fps}} = 4.8$.  An adaptive speed profile---linear from
\SI{6}{\metre\per\second} at \SI{20}{\metre} to \SI{10}{\metre\per\second} at
\SI{50}{\metre}---is used during search.

At every altitude in the operating range, the target appears in 11--17 frames
(Table~\ref{tab:frames}), providing substantial margin: a \emph{single}
detection above the 0.2 confidence threshold suffices to trigger
investigation.  The critical speed at which the target would appear in only
one frame is:
\begin{equation}\label{eq:critspeed}
  v_{\mathrm{crit}} = W_{\mathrm{gnd}} \cdot f_{\mathrm{fps}}
\end{equation}
which evaluates to \SI{66}{\metre\per\second} at \SI{20}{\metre} and
\SI{165}{\metre\per\second} at \SI{50}{\metre}---well beyond any multirotor
capability.  Frame coverage is therefore never the limiting factor.

\begin{table}[ht]
  \centering
  \caption{Frames on target at the adaptive speed profile.  Even at maximum
  speed the target is captured in $\geq$11 frames.}\label{tab:frames}
  \begin{tabular}{@{}c c c c c@{}}
    \toprule
    \textbf{Alt.\ (m)} & \textbf{Speed (m/s)} & \textbf{Frame gap (m)} &
      \textbf{Time in view (s)} & \textbf{Frames} \\
    \midrule
    20 &  6.0 & 1.25 & 2.28 & 11.0 \\
    30 &  7.3 & 1.53 & 2.82 & 13.5 \\
    35 &  8.0 & 1.67 & 3.00 & 14.4 \\
    50 & 10.0 & 2.08 & 3.43 & 16.5 \\
    \bottomrule
  \end{tabular}
\end{table}

\begin{figure}[ht]
  \centering
  \includegraphics[width=0.85\linewidth]{real_speed_vs_frames}
  \caption{Frames the target remains visible vs.\ drone speed at three
  altitudes (4.8\,FPS).  The green and red dashed lines mark the 3-frame
  smart-detect threshold and 1-frame minimum, respectively.}\label{fig:speed-frames}
\end{figure}

% ─────────────────────────────────────────────────────────────────────
\subsection{Coverage Efficiency}\label{sec:vp:coverage}

Search efficiency depends on both swath width and speed.  The coverage rate
(area surveyed per unit time) is:
\begin{equation}\label{eq:coverage}
  \dot{A} = H_{\mathrm{gnd}}(h) \times v(h)
\end{equation}
where $H_{\mathrm{gnd}}$ is the cross-track footprint and $v$ is the
altitude-adaptive speed.  At \SI{50}{\metre} the coverage rate is
approximately $343\;\text{m}^2/\text{s}$---roughly 1.7$\times$ that at
\SI{20}{\metre} ($82\;\text{m}^2/\text{s}$).  Starting the search at
\SI{50}{\metre} instead of \SI{35}{\metre} requires approximately 33\% fewer
scan lines and permits 25\% faster flight speed; the combined effect is that
the initial pass completes in roughly 60\% of the time.

The along-track overlap between consecutive frames at altitude~$h$ is:
\begin{equation}\label{eq:overlap}
  \eta = 1 - \frac{v(h)}{W_{\mathrm{gnd}}(h) \cdot f_{\mathrm{fps}}}
\end{equation}
At \SI{35}{\metre} and \SI{8}{\metre\per\second}, $\eta \approx 0.94$ (94\%
frame overlap).  Even at the maximum operating condition (\SI{50}{\metre},
\SI{10}{\metre\per\second}), overlap remains above 93\%.  No part of the
ground passes through the FOV without being imaged in multiple frames.

\begin{figure}[ht]
  \centering
  \includegraphics[width=0.85\linewidth]{real_coverage}
  \caption{Search efficiency vs.\ altitude.  The solid blue line shows
  coverage rate (m$^2$/s); the dashed pink line shows swath width.  Higher
  altitude improves efficiency at the cost of reduced target pixel
  size.}\label{fig:coverage}
\end{figure}

% ─────────────────────────────────────────────────────────────────────
\subsection{Combined Detection Envelope}\label{sec:vp:envelope}

The preceding analyses---pixel size, blur, and frame count---are
combined into a single scalar detection score to identify the operational
boundary:
\begin{equation}\label{eq:score}
  S(v, h) = S_{\mathrm{size}}(h) \times S_{\mathrm{frames}}(v, h)
    \times S_{\mathrm{blur}}(v, h)
\end{equation}
where each component is normalised to $[0, 1]$:
\begin{align}
  S_{\mathrm{size}}   &= \min\!\bigl(p_{\mathrm{model}}(h)\,/\,60,\;1\bigr)
    \label{eq:ssize} \\
  S_{\mathrm{frames}} &= \min\!\bigl(n_{\mathrm{frames}}(v,h)\,/\,5,\;1\bigr)
    \label{eq:sframes} \\
  S_{\mathrm{blur}}   &= 1 - \min\!\bigl(b_{\mathrm{px}}(v,h)\,/\,5,\;1\bigr)
    \label{eq:sblur}
\end{align}
A score of 60\,px in the model input and 5 frames on target represent ``fully
good'' conditions; 5\,px of blur represents total degradation.

Figure~\ref{fig:envelope} renders this score as a speed--altitude heatmap.
The green region ($S > 0.7$) covers the entire operating envelope
(6--\SI{10}{\metre\per\second}, 20--\SI{50}{\metre}), confirming that the
system operates well within its detection limits.  The $S = 0.5$ contour
(marginal detection) is only reached above \SI{55}{\metre} or at speeds
exceeding \SI{12}{\metre\per\second}---neither of which occurs during nominal
operation.

\begin{figure}[ht]
  \centering
  \includegraphics[width=0.88\linewidth]{real_detection_envelope}
  \caption{Combined detection score across the speed--altitude envelope.
  Green indicates high-confidence detection; the white contours mark
  $S = 0.3$, $0.5$, and $0.7$.  The operational regime (6--10\,m/s,
  20--50\,m) lies entirely within the $S > 0.7$ zone.}\label{fig:envelope}
\end{figure}

% ─────────────────────────────────────────────────────────────────────
\subsection{Summary of Limiting Factors}\label{sec:vp:summary}

Table~\ref{tab:limits} summarises the constraint analysis.  The dominant
operational limit is \emph{altitude} (pixel size), not speed or blur.  The
multi-pass descent strategy (50\,m $\to$ 40\,m $\to$ 32\,m) provides a
robust safety net: at 32\,m the dummy occupies 52\,px in the model
input---well above the detection floor.

\begin{table}[ht]
  \centering
  \caption{Summary of vision-system constraints within the operating
  envelope.}\label{tab:limits}
  \begin{tabular}{@{}l c p{6.5cm}@{}}
    \toprule
    \textbf{Factor} & \textbf{Limiting?} & \textbf{Reason} \\
    \midrule
    Motion blur     & No  & Global shutter + 1/1000\,s $\Rightarrow$ sub-pixel blur at all speeds \\
    Pixel size (${<}40$\,m)  & No  & Dummy height $>42$\,px in model input \\
    Pixel size (50\,m)       & Marginal & Dummy height 34\,px, width 7\,px \\
    Inference FPS   & No  & 4.8\,FPS gives 11--17 frames per flyover \\
    Frame gap       & No  & ${\leq}2.1$\,m gap at 10\,m/s; 93\%+ overlap \\
    Speed           & No  & Max 10\,m/s is far below 66\,m/s single-frame limit \\
    \bottomrule
  \end{tabular}
\end{table}

The speed profile could be increased to 15+\,m/s without impacting detection
performance; however, the resulting \SI{1.5}{\metre} GPS timing lag at
\SI{10}{\metre\per\second} (Section~\ref{sec:gps-accuracy}) makes speed
reduction a more effective strategy for improving localisation accuracy than
for preserving detection capability.
